\section{Conclusion}
\label{sec:conclusion}

In this paper, a neuro-symbolic architecture has been implemented and tested for the purpose of diagnosing breast cancer. DeepProbLog offers the advantage of using deep neural networks as predicates in a logic program whose predicates succeed with probabilistic rather than true/ false values in classic logic. This allows the user to trace the program and understands the reasons behind its final decision. The implemented architecture uses a pretrained residual network with 18 layers (ResNet18) that shares one trunk with four heads that estimate lesion kind, mass shape, mass margin, and calcification type in a mammogram image. The images are lesion
crops from the Curated Breast Imaging Subset of the Digital Database for Screening Mammography (CBIS-DDSM). The images annotations come from the Breast Imaging
Reporting and Data System (BI-RADS) lexicon. A second ResNet18 has been implemented as a black box network, that yields binary malignant/ non-malignant decisions, for the purpose of comparison.The paper describes the proposed neuro-symbolic approach as well as the training, validation, testing of both architectures and provides a discussion of the obtained results. One of the reasons that the obtained performance metrics are lower than expected, when using deep architectures, is that the resolution of the crops has been adjusted to 224 x 224 pixels due to hardware limitations (both memory size and training time). This corresponds to 70\% to 100\% of the original crop. \rev{The ROC-AUC difference is not significant.} Our focus was directed more towards explainability rather than towards performance. Explainability is helpful both for the designer of the architecture (it points to the head behind an error) as well as to the physician (it gives the evidence behind the diagnosis).
